\begingroup\small
\begin{longtable}{@{}P{.22\linewidth}P{.22\linewidth}P{.24\linewidth}P{.24\linewidth}@{}}
\caption{Integratieanalyse: bijdrage, grens en aangrenzende verantwoordelijkheid.}\label{tab:cross}\\
\toprule Instrument en niveau & Wel geadresseerd & Niet zelfstandig opgelost & Vereiste verbinding \\ \midrule\endfirsthead
\toprule Instrument en niveau & Wel geadresseerd & Niet zelfstandig opgelost & Vereiste verbinding \\ \midrule\endhead
ISO 704; terminologische methode \citep{S11} & Begripsafbakening, definities en aanduidingen. & Uitvoerbare constraints en distributie. & NL-SBB voor beschrijving; SKOS voor publicatie; mapping naar model. \\
NL-SBB; begrips\-beschrijvings\-model \citep{S14} & Consistente beschrijving van begrippen en bronnen. & Compleet domeinmodel of ontologische bewijsvoering. & SKOS-representatie, gezagsbron en MIM-koppeling; niet tot enkel SKOS-profiel reduceren. \\
SKOS; vocabulaire \citep{S31} & Labels, begripsstelsels, hiërarchie en mappings. & Objectidentiteit en volledige klasse-axiomatiek. & Expliciete relatie met OWL-klassen; geen automatische omzetting van broader naar subclass. \\
UFO/OntoUML; grondslagen en modelleertaal \citep{S40,S41,L09} & Analyse van categorieën, identiteit en afhankelijkheid. & Nederlandse normconformiteit of gegevensuitwisseling. & Selectieve toets van MIM-keuzen; OWL-vertaling vergt semantische beoordeling. \\
OntoClean; analysemethode \citep{L10} & Taxonomische keuzes via metaeigenschappen. & Productmetadata, operationeel beheer. & Alternatief of aanvulling op UFO-analyse van klassen. \\
MIM; metamodel \citep{S13} & Constructies voor conceptuele en logische informatiemodellen. & Rechtsgrond of volledige foundational ontology. & Begrippen naar objecttypen; attributen naar data-elementen; transformaties expliciteren. \\
NEN 3610; geo-informatienorm \citep{S17} & Gemeenschappelijke geo-modellerings\-grondslag binnen scope. & Alle niet-geografische masterdata of juridische identiteit. & Afstemming van geo-objectmodellering met MIM en domeincatalogus; volledige normtoets later. \\
ISO/IEC 11179-3, -31, -33, -6; metadata\-registratie \citep{S07,S08,S09,S10} & Gemeenschappelijke registratie, data-elementen, datasets en registratieproces. & Gehele productdienst, domeinontologie of alle beheerprocessen. & Conceptueel domein en waardedomein naar MIM en schema; datasetregistratie naar DCAT; beheer naar BOMOS. \\
ISO 8000-100/-110; uitwisseling masterdata \citep{S02,S03} & Scope voor karakteristieke data, syntax, semantic encoding en dataspecificatie. & Verplichte RDF-stack, algemene ontologie of interne MDM-architectuur. & Data-element en decodeerinformatie; implementatieprofiel apart normatief toetsen. \\
ISO 8000-120/-130/-140; kwaliteits\-informatie \citep{S04,S05,S06} & Provenance-, nauwkeurigheids- en volledigheidsinformatie binnen hun scope. & Automatische juistheid of vaste acceptatiedrempel. & PROV-O voor herkomstrepresentatie; DQV en casuscriteria voor metingen. \\
DAMA-DMBOK; kennisraamwerk \citep{S01} & Positionering van datamanagement; editie geverifieerd. & In dit dossier geen onderbouwing van specifieke boekvoorschriften. & Vervolglectuur voor verdieping; huidige requirements rusten op andere inhoudelijke bronnen. \\
RDF/RDFS; representatie en basissemantiek \citep{S28,S29} & Graaf, identifiers, klassen/eigenschappen en gevolgtrekking. & Waarheid, registratieplicht of gesloten datavalidatie. & OWL voor axioma’s; SHACL voor expliciete controles; provenance per bewering. \\
OWL; ontologietaal \citep{S30} & Formele axioma’s en logische gevolgtrekkingen. & Volledigheid van invoer of feitelijke juistheid. & Afzonderlijk SHACL-profiel; gemotiveerde mapping met begrippen en modellen. \\
SHACL; constraint\-specificatie \citep{S32} & Toetsbare constraints op de aangeboden graaf. & Volledige wereldkennis of juridische beoordeling. & Shapeversie aan model en gegevens koppelen; meetresultaat via kwaliteitsmetadata. \\
PROV-O; provenance-ontologie \citep{S33} & Entiteiten, activiteiten, agents en afleidingsrelaties. & Authenticiteit van de beschreven bron. & Bewijs en gezagsgrond apart; record-, model- en mappingversies refereren. \\
DQV; kwaliteits\-vocabulaire \citep{S34} & Representatie van dimensies, metrieken en metingen. & Concrete kwaliteitsnorm of volledige meetmethode. & ISO 25012/NORA voor classificatie; SHACL of externe toets als meetactiviteit. \\
ISO/IEC 25012 en NORA; kwaliteitsmodellen \citep{S12,S48} & Onderscheid tussen kwaliteitskenmerken en meetbare aspecten. & Casusgebonden drempels en één-op-één vocabulairemapping. & Metrieken expliciteren; NORA-kwaliteitsattribuut niet ongetoetst gelijkstellen aan DQV-concept. \\
OWL-Time; temporele ontologie \citep{S37} & Tijdstippen, intervallen en temporele relaties. & Automatisch valid/transaction-time-beheer. & Betekenis van elke tijdas en bronsysteem vastleggen; historie aan provenance koppelen. \\
URI/IRI en Cool URIs; identificatie en publicatierichtlijn \citep{S38,S39} & Syntax en onderscheid resource/document. & Persistentie zonder beheer of bewijs van co-identiteit. & Identifierbeleid, resolutie en versiebeheer; ontologische identiteit apart. \\
JSON-LD; representatie\-specificatie \citep{S35} & JSON verbinden met RDF-betekenis via contexten. & Domeindefinities of juiste contextselectie. & Contextversie aan dataset en API binden; rondreis controleren. \\
SPARQL; queryspecificaties \citep{S36} & Bevragen en protocollen voor RDF. & Productcontract, autorisatiebeleid of vereiste deployment. & Optionele toegang tot semantische graaf; vastleggen welke entailment geldt. \\
DCAT 3 / DCAT-AP / DCAT-AP-NL; catalogus en profielen \citep{S47,S45,S16} & Datasets, distributies, diensten en catalogusmetadata. & Volledig productcontract of betekenis van alle data-elementen. & Product over catalogusresources modelleren; profielversies en uitbreidingen afzonderlijk beheren. \\
FAIR; principes \citep{L11} & Vindbaarheid, toegankelijkheid, interoperabiliteit en herbruikbaarheid, ook voor machines. & Verplichte open toegang, specifieke technologie of gemeten AI-effect. & Identifier, metadata en provenance concretiseren per product en toegangsregime. \\
SEMIC; kernvocabularia \citep{S44} & Herbruikbare generieke begrippen voor Europese uitwisseling. & Juridische equivalentie met Nederlandse registratieobjecten. & Context- en versiegebonden mappings met domeinbegrippen. \\
OpenAPI; interface\-specificatie \citep{S19} & Beschrijving van HTTP-interfaces en hun structuren. & Volledige juridische en ontologische betekenis. & Semantische verwijzingen en contract aan interfaceversie koppelen. \\
NLGov REST API Design Rules; ontwerpregels \citep{S18} & Uniformere REST-interfaceafspraken binnen toepassingsgebied. & Gezag, ontologie of historische juistheid van inhoud. & OpenAPI-profiel, distributie en productrelease samen toetsen. \\
Digikoppeling; uitwissel\-afspraken \citep{S20} & Geprofileerde uitwisseling tussen systemen. & Betekenisgelijkheid tussen domeinmodellen. & Technisch uitwisselprofiel scheiden van semantisch contract. \\
FSC; federatieve uitwissel\-specificatie \citep{S21} & Deelnemers, contracten en technische vertrouwensrelaties. & Wettelijke authenticiteit of waarheid van een veld. & Afzenderidentiteit aan provenance koppelen zonder gezagsstatus over te nemen. \\
NL GOV CloudEvents; eventprofiel \citep{S22} & Gestandaardiseerde notificatie-envelop en verwijzingen. & Alle wijzigingsbetekenis, aflevering of bitemporale historie. & Object-/versiereferentie, changetype en herstelprotocol expliciet toevoegen binnen profiel. \\
BOMOS; beheerframework \citep{S23} & Organiseren van standaardbeheer en besluitvorming. & Volledig operationeel masterdataworkflowmodel. & Eigen lifecycle voor begrippen, modellen, producten en mappings vastleggen. \\
TOOI / MDTO; overheidssemantiek / bewaarmetadata \citep{S15,S24} & Actor- en informatieverwijzingen respectievelijk duurzaam toegankelijke informatieobjecten. & Complete masterdata-ontologie of universele bewaarplicht. & Waar relevant verantwoordelijke en bewaarde bewijsversie identificeren. \\
EIF / Interoperable Europe Act; kader / recht \citep{S43,S42} & Interoperabiliteits\-lagen respectievelijk toepasselijke juridische verplichtingen. & Eén voorgeschreven semantische stack. & Scopebesluit en interoperabiliteitsbeoordeling scheiden van technische validatie. \\
Stelseleisen, baseline, FDS; institutionele kaders \citep{S25,S26,S27} & Verantwoordelijk\-heden, kwaliteitszorg en federatieve afspraken. & Zelfstandig compleet semantisch metamodel. & Gezags-, feedback-, kwaliteits- en toepasselijkheidsmetadata per product. \\
\bottomrule
\end{longtable}
\endgroup